% Context from chapters/fourier.tex; for mathematical reference, not compilation.
\section{A Bit of Spectral Theory}

On a general domain $\XX$, the group-theoretic approach applies when $\XX=G$ is a group or when a group acts transitively on $\XX$. Another approach defines Fourier-like basis functions as eigenfunctions of a differential operator, following the characterization in Section~\ref{sec-fourier-pdes}. The Laplacian is particularly useful: it is second order, rotation invariant, and has natural counterparts on surfaces and graphs.


                                                
\subsection{On a Surface or a Manifold}

\begin{figure}[tbp]
\centering
\BookDrawing{tikz/fourier-laplacian-surface}
\caption{Computing the Laplacian on a surface.}
\label{fig-review-fourier-laplacian-surface}
\end{figure}
Let $\XX$ be a smooth compact connected Riemannian manifold of dimension $d$, without boundary. The Laplace--Beltrami operator can be characterized, for smooth $f$, by the small-ball mean-value expansion
\eq{
 (\Delta f)(x)=\lim_{\epsilon\to0}\frac{2(d+2)}{\epsilon^2}
 \left(\frac{1}{\mathrm{Vol}(B_\epsilon(x))}\int_{B_\epsilon(x)}f(y)\,\d\mu(y)-f(x)\right).
}
Here $\mu$ is Riemannian volume and $B_\epsilon(x)$ is the geodesic ball centered at $x$. The factor $\epsilon^{-2}$ is essential.

The operator $\Delta$ is unbounded and self-adjoint on its natural domain in $L^2(\XX)$; its resolvent is compact. It has a complete orthonormal family of eigenfunctions $(\phi_n)_{n\geq0}$, with eigenvalues, counted with multiplicity,
\eq{0=\lambda_0>\lambda_1\geq\lambda_2\geq\cdots\longrightarrow-\infty.}
The constant eigenfunction is $\phi_0=\mathrm{Vol}(\XX)^{-1/2}$, and the inner product is $\dotp{f}{g}_\XX=\int_\XX f(x)\overline{g(x)}\,\d\mu(x)$. Boundary conditions must be specified if the manifold has a boundary.

On the flat unit torus $\XX=(\RR/\ZZ)^d$,
\eq{
 \Delta f=\sum_{s=1}^d\frac{\partial^2f}{\partial x_s^2},
 \qquad\phi_n(x)=e^{2\pi\imath\dotp{n}{x}},
 \qquad\lambda_n=-4\pi^2\norm{n}^2,\quad n\in\ZZ^d.
}

                                                
\subsection{Spherical Harmonics}
\label{sec-spherical-harm}

\begin{figure}[tbp]
\centering
\BookDrawing{tikz/fourier-spherical-coordinates}
\caption{Spherical coordinates.}
\label{fig-review-fourier-spherical-coordinates}
\end{figure}
Consider the $(d-1)$-dimensional sphere $\SS^{d-1} = \enscond{x \in \RR^d}{\norm{x}_{\RR^d}=1}$.
  
The Laplacian has explicit eigenfunctions. For $d=3$, they are indexed by $n=(\ell,m)$:
\eq{
	\foralls \ell \in \NN, \quad
	\foralls m=-\ell,\ldots,\ell, \quad
	\phi_{\ell,m}(\th,\phi) = e^{\imath m \phi} P_{\ell}^m( \cos(\th) )
}
The corresponding eigenvalue is $\la_{\ell,m} = -\ell(\ell+1)$. Here $P_{\ell}^m$ are associated Legendre functions (polynomials multiplied by $(1-x^2)^{|m|